\documentclass[12pt,a4paper]{article}

\usepackage[utf8]{inputenc}
\usepackage[T1]{fontenc}
\usepackage{lmodern}
\usepackage[english]{babel}
\usepackage{amsmath, amssymb, amsthm}
\usepackage{tcolorbox}
\usepackage{enumitem}
\usepackage{array, booktabs}
\usepackage{multicol}
\usepackage[a4paper, margin=1in]{geometry}
\usepackage{fancyhdr}
\usepackage{hyperref}
\usepackage{cleveref}
\usepackage{graphicx}
\usepackage{caption}
\usepackage{footmisc}
\usepackage{xcolor}

\geometry{top=2.5cm, bottom=2.5cm, left=2.5cm, right=2.5cm, headheight=15pt}
\hypersetup{colorlinks=true, linkcolor=blue, filecolor=magenta, urlcolor=cyan, pdfpagemode=FullScreen}

\pagestyle{fancy}
\fancyhf{}
\lhead{\footnotesize \textit{Understanding Probabilistic Events and Information Theory in Mathematical Contexts}}
\rhead{\thepage}

% Custom colors for tcolorbox
\tcbuselibrary{theorems, skins, breakable}
\definecolor{thmcolor}{RGB}{200,230,255}
\definecolor{defcolor}{RGB}{230,250,220}
\definecolor{excolor}{RGB}{255,240,210}

% Theorem environment
\newtheorem{theorem}{Theorem}[section]
\newenvironment{thmbox}[1][]{%
  \begin{tcolorbox}[enhanced, breakable, colback=thmcolor, colframe=blue!50!black, title=#1, fonttitle=\bfseries, arc=3pt, boxrule=0.5pt]
}{\end{tcolorbox}}

% Definition environment
\newtheorem{definition}{Definition}[section]
\newenvironment{defbox}[1][]{%
  \begin{tcolorbox}[enhanced, breakable, colback=defcolor, colframe=green!50!black, title=#1, fonttitle=\bfseries, arc=3pt, boxrule=0.5pt]
}{\end{tcolorbox}}

% Lemma environment
\newtheorem{lemma}{Lemma}[section]
\newenvironment{lembox}[1][]{%
  \begin{tcolorbox}[enhanced, breakable, colback=yellow!10, colframe=orange!50!black, title=#1, fonttitle=\bfseries, arc=3pt, boxrule=0.5pt]
}{\end{tcolorbox}}

% Proposition environment
\newtheorem{proposition}{Proposition}[section]
\newenvironment{propbox}[1][]{%
  \begin{tcolorbox}[enhanced, breakable, colback=purple!5, colframe=purple!50!black, title=#1, fonttitle=\bfseries, arc=3pt, boxrule=0.5pt]
}{\end{tcolorbox}}

% Corollary environment
\newtheorem{corollary}{Corollary}[section]
\newenvironment{corbox}[1][]{%
  \begin{tcolorbox}[enhanced, breakable, colback=cyan!5, colframe=cyan!50!black, title=#1, fonttitle=\bfseries, arc=3pt, boxrule=0.5pt]
}{\end{tcolorbox}}

% Example environment
\newtheorem{example}{Example}[section]
\newenvironment{exbox}[1][]{%
  \begin{tcolorbox}[enhanced, breakable, colback=excolor, colframe=brown!50!black, title=#1, fonttitle=\bfseries, arc=3pt, boxrule=0.5pt]
}{\end{tcolorbox}}

% Remark environment
\newtheorem{remark}{Remark}[section]
\newenvironment{rembox}[1][]{%
  \begin{tcolorbox}[enhanced, breakable, colback=gray!5, colframe=gray!50, title=#1, fonttitle=\bfseries, arc=3pt, boxrule=0.5pt]
}{\end{tcolorbox}}

% Customize enumerate and itemize
\setlist[enumerate]{leftmargin=*, itemsep=2pt, topsep=4pt}
\setlist[itemize]{leftmargin=*, itemsep=2pt, topsep=4pt}

% Custom table environment with caption, placement, and column spec
\newenvironment{customtable}[3][htbp]{%
  % #1 = optional: placement (default: htbp)
  % #2 = mandatory: column specification (e.g., {ccc}, {l|c|r})
  % #3 = mandatory: table caption (goes above the table)
  \begin{table}[#1]
  \centering
  \renewcommand{\arraystretch}{1.2}
  \caption{#3}
  \begin{tabular}{#2}
}{%
  \end{tabular}
  \end{table}
}

\title{Understanding Probabilistic Events and Information Theory in Mathematical Contexts}

\begin{document}

\maketitle
\begin{abstract}
This lecture explores the relationship between probabilistic events and information theory, focusing on how different probabilistic events can yield varying amounts of information. The speaker discusses two independent probabilistic events and introduces the concept of joint information, asserting that it is the sum of individual information when events are independent. The logarithmic function is identified as the only function satisfying this relationship. The lecture also delves into the concept of average information in systems, emphasizing how observing a system can yield quantifiable information. Key examples illustrate how knowledge about a system can be measured, particularly in terms of uncertainty when no observations are made.
\end{abstract}

\section{Introduction}
This section introduces the exploration of probabilistic events and their relationship to information theory, focusing on the information derived from independent probabilistic events.

\section{Probabilistic Events}
In this section, we consider two probabilistic events, denoted as \(E_1\) and \(E_2\). The discussion revolves around which of these events provides more information.

\begin{defbox}[Probabilistic Events]
A probabilistic event is an occurrence that can be described by a probability measure, indicating the likelihood of its occurrence.
\end{defbox}

The speaker presents a scenario involving two professors, referred to as Professor L and Professor M, and discusses how their actions can lead to different amounts of information being conveyed. The key points include:

\begin{itemize}
    \item The first event involves Professor L interacting with students.
    \item The second event considers Professor M making choices about students without entering a room.
\end{itemize}

The speaker emphasizes that the second event is more informative because it is unexpected, leading to a greater amount of information being revealed.

\begin{remark}[Information and Probability]
The less probable an event is, the more information it conveys when it occurs. This is a fundamental concept in information theory.
\end{remark}

\section{Independent Probabilistic Events}
The discussion transitions to the concept of two independent probabilistic events. The speaker notes that when events are statistically independent, they can provide a combined measure of information.

\begin{defbox}[Statistical Independence]
Two events \(E_1\) and \(E_2\) are statistically independent if the occurrence of one does not affect the probability of the other occurring, mathematically expressed as:
\[
P(E_1 \cap E_2) = P(E_1) \cdot P(E_2).
\]
\end{defbox}

The implications of this independence are crucial for understanding how information is derived from multiple events.

\section{Joint Information and Logarithmic Functions}
This section discusses the concept of joint information in the context of independent probabilistic events and the significance of logarithmic functions in quantifying information.

\subsection{Joint Information}
When considering two independent probabilistic events, the joint information can be expressed as the sum of the individual information from each event. This relationship can be formally stated as follows:

\begin{thmbox}[Joint Information]
Let \( I(A) \) and \( I(B) \) represent the information from events \( A \) and \( B \), respectively. If \( A \) and \( B \) are independent, then the joint information \( I(A, B) \) is given by:
\[
I(A, B) = I(A) + I(B).
\]
\end{thmbox}

\subsection{Logarithmic Function}
The logarithmic function is identified as the unique function that satisfies the above relationship for joint information. Specifically, if we denote the information of an event as a function of its probability, the logarithm serves as the foundational mathematical tool for measuring information content.

\begin{defbox}[Logarithmic Function]
The logarithm, particularly in base 2, is commonly used in information theory to measure information in terms of bits. The relationship can be expressed as:
\[
I(X) = -\log_2(P(X))
\]
where \( P(X) \) is the probability of event \( X \).
\end{defbox}

\subsection{Average Information}
The concept of average information is introduced as a measure of the expected information that can be gained from observing a system. This average information quantifies the uncertainty associated with the system's state when no observations have been made.

\begin{rembox}[Average Information]
Average information provides insights into the behavior of a system over time and can be calculated using the probabilities of the system's various states.
\end{rembox}

\subsection{Information Received from Observations}
This section discusses how observing a system in a particular state provides information about that system. The key points include:

\begin{itemize}
    \item Upon observing the system, one receives a specific piece of information that quantifies the state of the system.
    \item The amount of information received can be measured, reflecting the knowledge gained from the observation.
\end{itemize}

\begin{defbox}[Quantity of Information]
The quantity of information received from observing a system is defined as the measure of information gained regarding the system's state.
\end{defbox}

\subsection{Uncertainty in Probabilistic Systems}
The discussion transitions to the concept of uncertainty in probabilistic systems. Key points include:

\begin{itemize}
    \item When no observations are made, the knowledge about the system is limited to the probabilities of its states, rendering the system somewhat mysterious.
    \item This lack of information leads to a measure of uncertainty about the system's behavior.
\end{itemize}

\begin{defbox}[Measure of Uncertainty]
The measure of uncertainty in a system is defined as the quantification of knowledge that is unknown due to the absence of observations.
\end{defbox}

\begin{remark}
Understanding the measure of uncertainty is crucial for analyzing probabilistic systems, as it provides insights into the potential information that can be gained through observation.
\end{remark}

\subsection{Minimum Probability and Information}
In this section, we analyze the minimum probability associated with a given probabilistic system. The key findings are:

\begin{itemize}
    \item The minimum value of the probability expression is determined to be zero.
    \item The probability for the first state, denoted as \( P(Z_1) \), is established to be equal to one, while all other terms are equal to zero.
\end{itemize}

This leads to the conclusion that if no prior knowledge about the system is available, the probability measure \( P \) can be considered as certain, yielding a value of zero for uncertainty. This reflects the principle that complete ignorance about a system results in a definitive measure of probability, which is \( P = 0 \).

\begin{defbox}[Principle of Probability]
If no information is known about a system, the measure of probability \( P \) is equal to zero, indicating complete uncertainty.
\end{defbox}

\begin{remark}
This principle emphasizes the relationship between knowledge and uncertainty in probabilistic systems, where the absence of information directly influences the probability measures assigned to various states.
\end{remark}

\section{References}
\begin{enumerate}
    \item Cover, T. M., & Thomas, J. A. (2006). \textit{Elements of Information Theory}. 2nd ed. Chapter 2: Information.
    \item Shannon, C. E. (1948). A Mathematical Theory of Communication. \textit{Bell System Technical Journal}, 27, 379-423.
    \item Kullback, S., & Leibler, R. A. (1951). On Information and Sufficiency. \textit{Annals of Mathematical Statistics}, 22(1), 79-86.
    \item MacKay, D. J. C. (2003). \textit{Information Theory, Inference, and Learning Algorithms}. Chapter 1: Information Theory.
    \item Jaynes, E. T. (1957). Information Theory and Statistical Mechanics. \textit{Physical Review}, 106(4), 620-630.
\end{enumerate}

\end{document}